\chapter{Connecting your laser}
\label{ch:connect}

Connecting a laser is three separate jobs, and it is worth keeping them separate
in your head: tell MeerK40t \emph{what} machine you own (a \emph{device}), tell it
\emph{how} to reach that machine (the \emph{interface}), and then prove that the
two agree by moving the head. This chapter walks all three. It ends with a
troubleshooting table you will come back to.

\section{Which driver is my machine?}

MeerK40t talks to a controller, not to a brand. Find your machine in the
catalogue below --- the entries are exactly what you see in the
\btn{Select Device} dialog of Figure~\ref{fig:devices-dialog} --- and note the
driver name in the right-hand column. The driver is what later chapters mean when
they say ``your device''.

\begin{center}\small
\begin{tabular}{@{}L{2.4cm}L{6.4cm}L{5.5cm}@{}}
\toprule
\textbf{Family} & \textbf{Entry in the list} & \textbf{Driver} \\
\midrule
K-Series CO\textsubscript{2} & \emph{K40-CO2-Laser (m2nano-Board) (Green/Blue)} &
Lihuiyu \\
& \emph{K40-CO2-Laser (m3nano-Board) (Purple/Blue)} & Lihuiyu \\
& \emph{K40 CO2 (GRBL-Controller)} & GRBL \\
& \emph{K40 CO2 --- MKS DLC32 (405$\times$285 mm)} & GRBL \\
& \emph{K50/K60-CO2-Laser (Ruida-Controller)} & Ruida \\
\midrule
Diode & \emph{Ortur Laser Master 2 (GRBL)},
\emph{Longer Ray5 (GRBL)}, \emph{Generic (GRBL-Controller)},
\emph{Diode-Laser (GRBL-Controller)} &
GRBL \\
\midrule
Other GRBL & \emph{GRBL-FluidNC (FluidNC-Controller)} & GRBL \\
\midrule
Moshi & \emph{CO2-Laser (Moshi-Board)} & Moshi \\
\midrule
Newly & \emph{RayLaser/U-SET}, \emph{Beijing SZTaiming}, \emph{Gama},
\emph{Artsign JSM-40U/3040U/3060U}, \emph{Rabbit40B}, \emph{Jinan Suke},
\emph{Lion laser} and a dozen more & Newly \\
\midrule
Galvo & \emph{Fibre-Laser (JCZ-Controller) (Non-MOPA)} or \emph{(MOPA)},
\emph{CO2 (JCZ-Controller)}, \emph{UV (JCZ-Controller)} & Balor \\
\bottomrule
\end{tabular}
\end{center}

\noindent If you cannot tell: a K40 with the original green or blue control board
is Lihuiyu; a diode engraver that speaks G-code is GRBL; a big
black-and-yellow CO\textsubscript{2} machine with a panel full of buttons is
Ruida; a fibre or UV marker with a scan-head and no gantry is Balor. When in
doubt, add a generic entry for the right family --- the defaults are only a
starting point.

\section{The device list}

Open \menu{Edit>Device-Manager}, or the \panel{Devices} pane
(\menu{Panes}), or \menu{Tools>Devices}. You get a table with one row per
configured machine:

\begin{center}\small
\begin{tabular}{@{}L{2.2cm}L{12.1cm}@{}}
\toprule
\textbf{Column} & \textbf{Meaning} \\
\midrule
\emph{Device} & The label you gave it. Double-click the cell to rename in
place. \\
\emph{Driver} & The internal driver class, for example \cmd{GRBLDevice} or
\cmd{LihuiyuDevice}. This is how you confirm which family a row belongs to. \\
\emph{Type} & The material source family: \cmd{Co2}, \cmd{Diode} and so on. It
decides which default operation settings are used. \\
\emph{Status} & The driver's own report --- \cmd{Idle}, \cmd{Running} or
\cmd{paused}. \\
\emph{Interface} & How the device is reached, for example \file{/dev/ttyUSB0:115200},
\file{192.168.1.5:23}, \file{usb} or \file{mock}. If this looks wrong, nothing
else will work. \\
\bottomrule
\end{tabular}
\end{center}

\noindent \textbf{The active device is drawn in red.} Only one device is active
at a time; that is the machine MeerK40t will drive. The buttons along the bottom
are \btn{Create New Device}, \btn{Duplicate}, \btn{Rename}, \btn{Remove},
\btn{Activate} and --- right-aligned, opening the settings --- \btn{Config}.
Grouping the rows is the whole point of having several: a shop with a K40 and a
diode machine keeps a device for each, duplicates one as a base for a second
identical machine, and switches with \btn{Activate} or by double-clicking the
row. The active device cannot be removed, and \btn{Remove} asks for confirmation
because it cannot be undone.

\begin{tipbox}[The same thing from the console]
\begin{lstlisting}[style=konsole]
device                      # list the registered device entries
device add grbl-generic -l "My diode"
device activate "My diode"
device duplicate "My diode" "Spare diode"
device delete "Spare diode"
devinfo                     # current logical and native position
\end{lstlisting}
These commands are what the dialog is doing behind the scenes, and they are the
quickest way to script a machine setup.
\end{tipbox}

\section{Telling MeerK40t how big the bed is}

Press \btn{Config} on the active row. Every driver has a Configuration window
whose first tab contains the geometry settings. Get these right before you draw
anything, because the bed size defines the scene's bed rectangle and the
origin --- and therefore where your job will be cut.

\begin{center}\small
\begin{tabular}{@{}L{3.6cm}L{2.9cm}L{7.8cm}@{}}
\toprule
\textbf{Label} & \textbf{Typical default} & \textbf{What to do} \\
\midrule
\btn{Width}, \btn{Height} & 235\unit{mm} (GRBL), 310$\times$210\unit{mm} (Lihuiyu,
Newly), 330$\times$210\unit{mm} (Moshi), 24$\times$16\unit{in} (Ruida),
110\unit{mm} lens (Balor) & Measure your machine's real travel and enter it. The
specification is often optimistic. \\
\btn{Flip X}, \btn{Flip Y} & off, except GRBL Y and Balor X/Y & These make the
on-screen axes match the machine. GRBL devices usually want \emph{Flip Y} on so
that Y grows into the bed rather than away from it. \\
\btn{Swap X and Y} & off (on for Balor) & For machines whose axes are mounted
rotated. Applied before the flips. \\
\btn{Force Declared Home} & \cmd{auto} & Which corner is the machine's zero:
\emph{auto}, \emph{top-left}, \emph{top-right}, \emph{bottom-left},
\emph{bottom-right} or \emph{center}. \\
\btn{X-Axis} / \btn{Y-Axis} scale & 1.000 & Corrections for machines that move
the wrong distance --- calibrate with a measured test burn, not by guesswork. \\
\btn{Laserspot} & 0.3\unit{mm} & Your focal spot size. Used by the raster
planner and by ``consider laserdot''. \\
\bottomrule
\end{tabular}
\end{center}

\begin{notebox}[Two settings you will see but should ignore]
\emph{X-Margin} and \emph{Y-Margin} exist in several drivers' settings, with a
default of \cmd{0}, but they are hidden from the GUI and marked in the program's
own source as not working yet. Do not build a workflow on them. Position jobs on
the scene instead (Chapter~\ref{ch:where}).
\end{notebox}

\section{Connecting each family}

\subsection*{Lihuiyu / K40 over USB}

The K40's controller board is reached through a CH341 USB-to-serial chip. In the
\tabname{Interface} tab, choose the radio button \btn{USB} (or \btn{Networked},
see below, or \btn{Mock} for a pretend machine). The USB settings let you pin
down which board you mean when several are attached: \btn{Restrict Multiple
Lasers} offers a CH341 version, a device index, and a serial-number match. The
\btn{Write Buffer} group limits how much is queued at once (default 1500) ---
relevant on the oldest boards.

Two board-level settings on the first tab matter: \btn{Board} (\cmd{M2} is the
classic green board, \cmd{M3} the newer one) and \btn{Maximum Laser strength},
which caps output as a percentage of the tube's rating. The default cap is
deliberately low and the program warns when you raise it: setting it too high can
damage a tube that is smaller than the software assumes.

\subsection*{GRBL}

Pick the interface in the \tabname{Interface} tab with \btn{Interface Type}:

\begin{center}\small
\begin{tabular}{@{}L{2.6cm}L{3.4cm}L{8.3cm}@{}}
\toprule
\textbf{Interface} & \textbf{Settings} & \textbf{Notes} \\
\midrule
\emph{Serial} & \btn{Serial Port} + \btn{Baud Rate} (115200 default) & The
normal USB cable case. The port list is built from the ports the system reports;
on Linux expect \file{/dev/ttyUSB0} or \file{/dev/ttyACM0}. \\
\emph{TCP-Network} & \btn{Address} + \btn{Port} (23) & For a networked or
Wi-Fi-equipped controller; also used by the K40 MKS DLC32 profile, which defaults
to \file{192.168.10.90:8080}. \\
\emph{WebSocket-Network} & \btn{Address} + \btn{Port} (81) & For FluidNC-style
firmware. \\
\emph{Mock} & --- & A fake machine that accepts everything. Excellent for
practising the workflow without hardware. \\
\bottomrule
\end{tabular}
\end{center}

\noindent Four other GRBL settings save a lot of time:

\begin{itemize}
  \item \btn{Require Validator} with \btn{Welcome Validator} (\cmd{Grbl}) ---
        refuses to send anything until the controller has introduced itself. Turn
        it on once your connection is reliable; it turns ``nothing happens'' into
        an explicit error.
  \item \btn{Reset on connect} --- sends a soft reset when the port opens. Useful
        when a controller wakes up confused.
  \item \btn{Reset on 'Clear Alarm'} (\setting{extended_alarm_clear}) --- also
        resets the controller when you clear an alarm, for machines that need it.
  \item \btn{Device has endstops} --- enables real homing. With endstops on a
        top-left-homing board, also enable \btn{Sequential homing}, which homes Y
        before X; without it such boards stop short by around 50\unit{mm}.
\end{itemize}

\noindent The ribbon for a GRBL device adds \btn{Clear Alarm}, \btn{Home}
(right-click for physical home), \btn{Red Dot On}/\btn{Red Dot Off} for
machines with a pilot laser, and --- if ESP3D upload is enabled --- the ESP3D
upload and control buttons.

\subsection*{Ruida}

Ruida machines are configured as either \btn{USB} (a serial interface field
accepting \file{/dev/...} or \cmd{COM\emph{n}}) or \btn{Networked}, which is UDP
with an address setting. Two facts to know before you rely on it: the driver is
marked in the project's own documentation as under revision, and users of at
least one tested controller report that \emph{Flip X} and a raster
\emph{Overscan} of zero are required. Test on scrap.

MeerK40t can also \emph{pretend} to be a Ruida controller --- useful when you want
to drive a K40 from software that only speaks Ruida
(\cmd{ruidacontrol}, Chapter~\ref{ch:remote}).

\subsection*{Moshi and Newly}

Both use USB, both are older CO\textsubscript{2} machines, and both are
reported as stable. Moshi boards talk through the same CH341 chip as the K40, so
the USB-restriction settings look familiar. Newly adds a machine-specific
concept: the controller stores the job in one of ten numbered \emph{files}
(\btn{Output File}, default 0). File~0 executes immediately; the others can be
sent and then replayed from the machine's own panel, which is what the ribbon
buttons \btn{File 0}\dots\btn{File 9}, \btn{Send Only}, \btn{Draw Frame},
\btn{Move Frame} and \btn{Replay} are for.

\subsection*{Balor / JCZ galvo}

Galvo machines connect over USB only --- there is no host and port to set. What
matters in the configuration is the optics and the machine index:
\btn{Lens Size} (default 110\unit{mm}) describes your field, and enabling a
correction file (\btn{Enable} + \btn{File}, a \file{.cor} file from the lens
supplier) makes the marked geometry match the design. \btn{Machine index to
select} picks the controller when more than one is attached; \btn{Check serial
no} lets you match a specific unit. The \tabname{Timings} tab holds the
laser on/off and jump delays --- change these only with a documented reason, since
they control whether corners are burned through or left open.

\section{USB permissions on Linux}

If the machine is not found on Linux, or only works when you run the program as
root, the cause is almost always permission on the USB device. Create
\file{/etc/udev/rules.d/99-meerk40t.rules} with the vendor rules for the chips
MeerK40t uses:

\begin{lstlisting}[style=file]
# CH341 USB-to-serial chips (used by Lihuiyu laser controllers)
SUBSYSTEM=="usb", ATTR{idVendor}=="1a86", ATTR{idProduct}=="*", MODE="0666"
# FTDI USB-to-serial chips (used by Ruida laser controllers)
SUBSYSTEM=="usb", ATTR{idVendor}=="0403", ATTR{idProduct}=="*", MODE="0666"
# BJJCZ galvo controllers
SUBSYSTEM=="usb", ATTR{idVendor}=="9588", ATTR{idProduct}=="*", MODE="0666"
# GSI Group galvo controllers
SUBSYSTEM=="usb", ATTR{idVendor}=="28e9", ATTR{idProduct}=="*", MODE="0666"
# Philips chips (used by Newly controllers)
SUBSYSTEM=="usb", ATTR{idVendor}=="0471", ATTR{idProduct}=="*", MODE="0666"
\end{lstlisting}

\noindent then apply them and add yourself to the serial group:

\begin{lstlisting}[style=konsole]
sudo udevadm control --reload-rules && sudo udevadm trigger
sudo usermod -a -G dialout $USER
\end{lstlisting}

\noindent Log out and back in for the group change. Do not run MeerK40t with
\cmd{sudo}: it is not recommended, and it produces configuration files your
normal user cannot then read.

\section{Proving the connection works}

Before you cut anything, answer one question: can MeerK40t talk to the machine?
Three windows tell you, all of them under \menu{Tools}:

\begin{itemize}
  \item \btn{Controller} (per device) --- the driver's own window, with
        connection state, packet counters, a \btn{Reset statistics} button,
        \btn{Clear Status-Log} and the raw traffic. If the connection is alive,
        the counters climb.
  \item \btn{Buffer} (\panel{BufferView}) --- shows the command flow queued to
        the device. Useful when a job appears to hang: is the queue empty because
        the job finished, or because it never started?
  \item \panel{UsbConnect} --- tails the raw USB conversation and lets you type
        into the pipe by hand. This is the deepest view available.
\end{itemize}

\noindent In the console, \cmd{devinfo} prints the current logical and native
position, and \cmd{device} lists your devices. If the interface column in the
device list shows a plausible port and the controller window shows traffic, the
link is up.

\section{Moving the head safely}
\label{sec:jog}

The \panel{Jog} pane is where you learn to trust the machine. It has nine
direction buttons plus \btn{Home}, \btn{Unlock rail}, \btn{Confine} and \btn{Lock
rail}.

\begin{machinebox}{Do your first movements with the beam off}
First moves are the ones where the head tries to leave the machine. Keep the
enclosure closed, keep your hand near the emergency stop, and do not have
material --- or your fingers --- under the head. If a direction button sends the
head the wrong way, do not chase it with software: stop, and fix \emph{Flip X},
\emph{Flip Y} or \emph{Swap X and Y} in the device configuration.
\end{machinebox}

\noindent Three distinctions worth learning on day one:

\begin{description}[leftmargin=1.4em,style=nextline]
  \item[Home versus physical home] Left-clicking \btn{Home} sends the head to the
        \emph{defined} home --- the origin the driver and your settings describe.
        Right-clicking it runs the \emph{physical} home, which seeks the
        controller's endstops. On GRBL the physical home sends \cmd{\$H} (or
        \cmd{\$HY} then \cmd{\$HX} when sequential homing is on). Machines without
        endstops simply move to 0,0.
  \item[Locked versus unlocked rail] \btn{Unlock rail} lets the head be moved by
        hand; \btn{Lock rail} re-engages it. On some drivers the rail locks itself
        automatically after a job (\setting{autolock}).
  \item[Confined movement] \btn{Confine} limits jogging to the bed size. Turning
        confinement off lets you drive the head outside the declared area ---
        which on GRBL means soft-limit alarms, and on anything else means
        crashing into the frame. The program warns you when you switch it off.
\end{description}

\noindent Jog distance is set in the \panel{Distances} pane
(\panel{Jog Distance:}); the default is 10\unit{mm} per press, and holding a
button repeats the move. \btn{Move Laser to:} moves to typed X/Y coordinates, and
the nine buttons beside it are saved positions --- left-click to go there,
right-click to store the current position. They persist between sessions.

\section{When it does not connect}

\begin{center}\small
\begin{tabular}{@{}L{6.1cm}L{8.2cm}@{}}
\toprule
\textbf{Symptom} & \textbf{Check, in this order} \\
\midrule
No device entry matches your machine & Add a generic entry for the family
(Chapter~\ref{ch:connect}); the defaults are editable. \\
Interface column says \cmd{UNCONFIGURED} or is empty & Choose the serial port (or
type the network address) in the Interface tab; replug the cable and reopen the
list. \\
Port exists but nothing moves & Baud rate (115200 on almost all GRBL boards),
\emph{require validator} and the welcome string, and \emph{reset on connect}. \\
Works only with \cmd{sudo} (Linux) & Install the udev rules above and add yourself
to \cmd{dialout}. \\
GRBL answers with an alarm & \btn{Clear Alarm}; if the alarm returns
immediately, the machine is hitting soft limits --- check bed size, flips and
homing corner. \\
Head moves the wrong way & \emph{Flip X}/\emph{Flip Y}/\emph{Swap XY}. Change one
at a time and re-test. \\
Head moves the wrong distance & Axis scale factors; calibrate with a measured
move rather than arithmetic. \\
Job stops halfway with the queue empty & Check the controller window for an alarm
or an error, and the console for the last message received. \\
Two K40s, only one responds & \emph{Restrict Multiple Lasers}: set the device
index or serial number for each. \\
\bottomrule
\end{tabular}
\end{center}

\begin{tryit}{Commission a machine in ten minutes}
Work through this once, on any machine, and keep the results in your notes.
\begin{enumerate}[label=\arabic*.,leftmargin=1.4em]
  \item Add or activate the correct device. Confirm the \emph{Driver} and
        \emph{Type} columns look right.
  \item Enter the real bed size, measured with a tape, and set \emph{Force
        Declared Home} to the corner your machine actually uses.
  \item Open the \btn{Controller} window and confirm traffic. Note the interface
        string.
  \item In the \panel{Jog} pane, set the jog distance to 10\unit{mm} and jog in
        all four directions, watching the head.
  \item Home the machine (left-click), then right-click to run the physical home
        if the machine has endstops.
  \item Jog to a known point, then use \btn{Move Laser to:} to send the head to
        X~0, Y~0 and confirm it lands where the machine thinks zero is.
  \item Lock the rail, then unlock it and nudge the head by hand and back.
  \item Save the console transcript or take a screenshot of the device
        configuration. When something misbehaves in three months, this is your
        reference.
\end{enumerate}
\end{tryit}

\noindent Full per-driver settings tables, including every default mentioned in
this chapter, are collected in Appendix~\ref{app:drivers}.
